\documentclass[11pt,a4paper]{article}

\usepackage[utf8]{inputenc}
\usepackage[T1]{fontenc}
\usepackage{lmodern}
\usepackage[margin=0.78in]{geometry}
\usepackage[table]{xcolor}
\usepackage{array}
\usepackage{tabularx}
\usepackage{booktabs}
\usepackage{enumitem}
\usepackage{titlesec}
\usepackage{fancyhdr}
\usepackage{lastpage}
\usepackage{hyperref}
\usepackage[normalem]{ulem}

\definecolor{TextInk}{HTML}{233142}
\definecolor{HeadingBlue}{HTML}{1D3557}
\definecolor{AccentTeal}{HTML}{2A7F92}
\definecolor{AccentGold}{HTML}{C58B4E}
\definecolor{RuleGray}{HTML}{D6E0E5}
\definecolor{TableStripe}{HTML}{F3F8FA}
\definecolor{SoftTint}{HTML}{F7F4EE}
\definecolor{SoftBlue}{HTML}{EEF6F8}
\definecolor{SoftGreen}{HTML}{EEF7F0}
\definecolor{SoftRed}{HTML}{FBF0EF}

\hypersetup{
  colorlinks=true,
  linkcolor=HeadingBlue,
  urlcolor=AccentTeal,
  pdftitle={IB Geography SL Case Study Bank Day - 2026-05-01},
  pdfauthor={IB Geography SL Preparation}
}

\color{TextInk}
\setlength{\parskip}{0.5em}
\setlength{\parindent}{0pt}
\setlength{\tabcolsep}{5pt}
\renewcommand{\arraystretch}{1.19}
\setlist[itemize]{leftmargin=1.35em,itemsep=3pt,topsep=4pt}
\setlist[enumerate]{leftmargin=1.5em,itemsep=3pt,topsep=4pt}

\newcolumntype{Y}{>{\raggedright\arraybackslash}X}
\newcolumntype{L}[1]{>{\raggedright\arraybackslash}p{#1}}

\titleformat{\section}
  {\normalfont\huge\sffamily\bfseries\color{HeadingBlue}}
  {}{0pt}{}
  [\vspace{0.12em}{\color{AccentGold}\titlerule[1pt]}]

\titleformat{\subsection}
  {\normalfont\Large\sffamily\bfseries\color{AccentTeal}}
  {}{0pt}{}

\titleformat{\subsubsection}
  {\normalfont\large\sffamily\bfseries\color{HeadingBlue}}
  {}{0pt}{}

\titlespacing*{\section}{0pt}{1.25em}{0.55em}
\titlespacing*{\subsection}{0pt}{0.9em}{0.25em}
\titlespacing*{\subsubsection}{0pt}{0.7em}{0.2em}

\pagestyle{fancy}
\setlength{\headheight}{14pt}
\fancyhf{}
\fancyhead[L]{\sffamily\color{AccentTeal}Fri 2026-05-01}
\fancyhead[R]{\sffamily\color{HeadingBlue}IB Geography SL}
\fancyfoot[C]{\sffamily\color{HeadingBlue}\thepage\ of \pageref{LastPage}}
\renewcommand{\headrulewidth}{0.4pt}
\renewcommand{\footrulewidth}{0pt}

\newcommand{\term}[1]{\textcolor{HeadingBlue}{\textbf{#1}}}
\newcommand{\exam}[1]{\textcolor{AccentTeal}{\textbf{#1}}}
\newcommand{\boxnote}[3]{%
\noindent\colorbox{#1}{%
  \begin{minipage}{0.965\textwidth}
  \sffamily
  \textbf{\textcolor{HeadingBlue}{#2}} #3
  \end{minipage}
}}

\begin{document}

\begin{center}
  {\sffamily\bfseries\color{HeadingBlue}\LARGE IB Geography SL Study Pack}\\[0.25em]
  {\sffamily\bfseries\color{AccentTeal}\Large Friday 2026-05-01}\\[0.35em]
  {\sffamily\color{TextInk}Case-Study Bank Day: Urban Environments + Resource Stewardship}
\end{center}

\vspace{0.4em}

\boxnote{SoftBlue}{How this pack follows the plan.}{The May 1 schedule is a
weekday case-study bank session. \term{Audio 1} covers \term{Urban
environments} named examples and statistics. \term{Audio 2} covers
\term{Global resource consumption and security} case studies and stewardship
examples. The \term{25 minute} block creates a rapid-recall sheet using:
\term{place} $\rightarrow$ \term{issue} $\rightarrow$ \term{causes}
$\rightarrow$ \term{impacts} $\rightarrow$ \term{response} $\rightarrow$
\term{evaluation}.}

\section{Today's Target}

By the end of this session, you should have:

\begin{itemize}
  \item one exam-ready \term{Urban environments} class case: Detroit
  gentrification and uneven regeneration
  \item one optional urban management comparison case: Curitiba BRT and
  transport planning
  \item one exam-ready \term{Resources} class case: China waste-import
  restrictions and e-waste / circular economy
  \item one resource stewardship case: Singapore water security and the Four
  National Taps
  \item a printable rapid-recall sheet with the six-link case-study chain
  \item no invented statistics: any uncertain detail is flagged for checking
\end{itemize}

\subsection{Session Structure}

\rowcolors{2}{TableStripe}{white}
\begin{tabularx}{\textwidth}{L{0.18\textwidth}Y L{0.28\textwidth}}
\toprule
\textbf{Part} & \textbf{Focus} & \textbf{Output} \\
\midrule
Audio 1 & Urban named examples and statistics: Detroit first, Curitiba as a
comparison if useful & Two urban case chains \\
Audio 2 & Resources case studies: China e-waste / waste imports and Singapore
water stewardship & Two resource case chains \\
25 min practice & Rapid-recall sheet with place, issue, causes, impacts,
response, evaluation & One completed case-bank sheet \\
\bottomrule
\end{tabularx}
\rowcolors{2}{}{}

\boxnote{SoftTint}{Case-study discipline.}{Use the examples you have actually
studied first. The student case-study sheet currently lists
\term{Gentrification in Detroit} for Urban environments and \term{China no
longer accepting our e-waste} for Resources. Curitiba and Singapore are
included as compact comparison or stewardship examples; mark them as
\term{backup} unless they match your class notes.}

\section{Case-Study Rule}

\subsection{The Six-Link Chain}

Every useful case study should fit this sequence:

\begin{center}
{\sffamily\bfseries\color{HeadingBlue}
Place $\rightarrow$ Issue $\rightarrow$ Causes $\rightarrow$ Impacts
$\rightarrow$ Response $\rightarrow$ Evaluation}
\end{center}

\rowcolors{2}{TableStripe}{white}
\begin{tabularx}{\textwidth}{L{0.18\textwidth}Y Y}
\toprule
\textbf{Link} & \textbf{What to store} & \textbf{Exam warning} \\
\midrule
Place & City, country, scale, and date if useful & Avoid vague examples such
as ``a city in Brazil'' \\
Issue & The exact geography problem or process & Do not turn every case into a
general story \\
Causes & Physical, social, economic, political, or demographic drivers & More
than one cause usually makes the answer stronger \\
Impacts & Social, economic, environmental, and spatial effects & Identify who
is affected and where \\
Response & Policy, planning, management, stakeholder action, or technology &
Describe the response briefly, then judge it \\
Evaluation & Success, limitation, uneven outcome, cost, equity, or side effect
& This is where many marks are gained \\
\bottomrule
\end{tabularx}
\rowcolors{2}{}{}

\section{Audio 1 Notes}

\subsection{Urban Environments Named Examples}

\subsubsection{Case 1: Detroit, United States}

\rowcolors{2}{TableStripe}{white}
\begin{tabularx}{\textwidth}{L{0.18\textwidth}Y}
\toprule
\textbf{Case link} & \textbf{Exam-ready content} \\
\midrule
Place & Detroit, Michigan, United States \\
Issue & Gentrification and uneven regeneration after deindustrialization,
population loss, vacancy, and fiscal stress \\
Statistics & U.S. Census: Detroit had \term{1,849,568} people in 1950;
\term{639,111} in the 2020 Census; \term{645,705} estimated on July 1, 2024.
QuickFacts also gives \term{32.7\%} poverty and median household income of
about \term{\$39,938} for 2020--2024. \\
Causes & Automobile-industry restructuring, suburbanization, population loss,
racial and income segregation, vacancy, and later selective reinvestment \\
Impacts & Improved image and activity in reinvested districts, but uneven
benefits, rent pressure or exclusion risk, and persistent poverty / service
gaps in other areas \\
Response & Regeneration, reuse of buildings, targeted investment, possible
housing and neighborhood policy from class notes \\
Evaluation & Regeneration is only fully sustainable when lower-income
residents benefit, not just investors, commuters, tourists, or new
higher-income residents \\
\bottomrule
\end{tabularx}
\rowcolors{2}{}{}

\boxnote{SoftGreen}{Best use.}{Detroit is strongest for
\term{deindustrialization}, \term{urban decline}, \term{gentrification},
\term{regeneration}, \term{social stress}, and \term{uneven sustainable urban
management}. It is not automatically the best case for transport or
environmental management unless your class notes provide those details.}

\subsubsection{Detroit Model Paragraph}

Detroit shows how post-industrial urban change can create both decline and
selective regeneration. The city's population fell from about \term{1.85
million} in 1950 to \term{639,111} in the 2020 Census, reflecting long-term
deindustrialization, suburbanization, and population loss. This reduced the tax
base, increased vacancy, and contributed to uneven services. Later
reinvestment and gentrification in selected districts can improve image,
services, and economic activity, but the strategy is only partly sustainable
where poverty remains high and lower-income residents do not gain affordable
housing, access to jobs, or improved neighborhood services. Therefore Detroit
is best used as a case of \term{uneven regeneration}, not as a simple success
story.

\subsubsection{Case 2: Curitiba, Brazil}

\rowcolors{2}{TableStripe}{white}
\begin{tabularx}{\textwidth}{L{0.18\textwidth}Y}
\toprule
\textbf{Case link} & \textbf{Exam-ready content} \\
\midrule
Place & Curitiba, Parana, Brazil \\
Issue & Sustainable transport, BRT, and urban growth management \\
Statistics & New Development Bank: Curitiba municipality has an estimated
population of \term{1.9 million}; NDB approved \term{US\$75 million} in 2020 for BRT
rideability improvements; the East-West corridor component includes \term{33}
stations over \term{22.5 km}; the South Corridor includes \term{13} stations
over \term{7.5 km}. \\
Causes & Urban growth, congestion pressure, need for efficient public
transport, and need for a lower-cost alternative to heavy rail \\
Impacts & BRT can improve mobility, reduce car dependence, and help structure
urban development around corridors \\
Response & Bus rapid transit, dedicated corridors, feeder-network integration,
station and accessibility improvements, traffic control, pavement upgrades \\
Evaluation & BRT can be flexible and lower-cost, but long-term success depends
on speed, maintenance, funding, accessibility, ridership, and whether peripheral
users are well served \\
\bottomrule
\end{tabularx}
\rowcolors{2}{}{}

\boxnote{SoftTint}{Use carefully.}{Curitiba is a backup comparison case unless
it appears in your class notes. Its strongest value is for sustainable
transport and future urban systems. Pair it with Detroit when you need to
compare \term{regeneration} with \term{planned transport management}.}

\section{Audio 2 Notes}

\subsection{Resources Case Studies And Stewardship}

\subsubsection{Case 3: China, E-Waste, And Waste Imports}

\rowcolors{2}{TableStripe}{white}
\begin{tabularx}{\textwidth}{L{0.18\textwidth}Y}
\toprule
\textbf{Case link} & \textbf{Exam-ready content} \\
\midrule
Place & China plus waste-exporting countries in high-consumption economies \\
Issue & Imported solid waste / e-waste, circular economy, and externalized
pollution risk \\
Statistics & China banned \term{24 types} of solid waste in four classes at
the start of 2018, including unsorted wastepaper and scrap plastic. The Global
E-waste Monitor 2024 reports \term{62 billion kg} of e-waste generated
globally in 2022, with \term{22.3\%} formally collected and recycled. \\
Causes & High consumption, short product life cycles, weak repair systems,
cheap waste trade, inadequate domestic recycling capacity, and demand for raw
materials \\
Impacts & Pollution and health risk in receiving areas, lost valuable
materials, disruption for exporters after the ban, stockpiling, landfill or
incineration pressure \\
Response & Import restrictions, stricter permits, circular economy, producer
responsibility, repair, reuse, collection systems, and domestic processing \\
Evaluation & China's policy reduced one pathway of foreign waste pressure, but
it did not solve global waste generation; stewardship requires source
reduction, safer recycling, and product design changes \\
\bottomrule
\end{tabularx}
\rowcolors{2}{}{}

\subsubsection{Case 4: Singapore Water Stewardship}

\rowcolors{2}{TableStripe}{white}
\begin{tabularx}{\textwidth}{L{0.18\textwidth}Y}
\toprule
\textbf{Case link} & \textbf{Exam-ready content} \\
\midrule
Place & Singapore \\
Issue & Water security in a small, land-scarce, highly urbanized country \\
Statistics & Four National Taps: local catchment, imported water, NEWater,
desalination. About \term{two-thirds} of land is water catchment; over
\term{8,000 km} of drains and canals move rainwater to \term{17 reservoirs}.
Imported water from Johor can be up to \term{250 million gallons per day}
under the 1962 agreement. Singapore has \term{five} NEWater factories and
\term{five} desalination plants. Demand is about \term{440 million gallons per
day} and is expected to almost double by 2065. \\
Causes & Limited natural freshwater, limited land, high population and
industrial demand, imported-water dependence, and climate uncertainty \\
Impacts & Risk of water insecurity, pressure on infrastructure, demand growth,
and need for resilient supply during dry periods \\
Response & Four National Taps, NEWater, desalination, catchment expansion,
water pricing, efficiency labels, industry regulation, conservation education \\
Evaluation & Singapore shows that technology and governance can manage
scarcity, but the system is costly, partly energy intensive, and depends on
public trust, technical capacity, demand management, and climate resilience \\
\bottomrule
\end{tabularx}
\rowcolors{2}{}{}

\section{Comparison Toolkits}

\subsection{Urban Comparison}

\rowcolors{2}{TableStripe}{white}
\begin{tabularx}{\textwidth}{L{0.18\textwidth}Y Y}
\toprule
\textbf{Criterion} & \textbf{Detroit} & \textbf{Curitiba} \\
\midrule
Main process & Deindustrialization, urban decline, regeneration,
gentrification & Planned public transport and sustainable urban management \\
Best prompts & Gentrification, urban social stress, uneven regeneration,
deindustrialization & Sustainable transport, future cities, congestion,
environmental management \\
Core statistic & 1.85 million people in 1950; 639,111 in 2020 & 1.9 million
people; NDB-funded BRT upgrades of US\$75 million \\
Evaluation angle & Recovery is spatially uneven and must be judged by
affordability and inclusion & BRT is useful but needs continuing investment,
ridership, speed, and accessibility \\
\bottomrule
\end{tabularx}
\rowcolors{2}{}{}

\subsection{Resource Comparison}

\rowcolors{2}{TableStripe}{white}
\begin{tabularx}{\textwidth}{L{0.18\textwidth}Y Y}
\toprule
\textbf{Criterion} & \textbf{China / e-waste} & \textbf{Singapore / water} \\
\midrule
Resource issue & Waste, valuable materials, hazardous substances, circular
economy & Water security, scarcity, demand growth, technological stewardship \\
Main driver & High consumption and short product life cycles in a globalized
waste system & Limited natural freshwater plus high urban and industrial demand \\
Response & Import restrictions, circular economy, safer recycling, producer
responsibility & Four National Taps, NEWater, desalination, conservation,
pricing, regulation \\
Evaluation angle & Policy shifts waste geography unless source reduction and
safe recycling improve & Diversification improves security but has energy,
cost, dependency, and demand-management limits \\
\bottomrule
\end{tabularx}
\rowcolors{2}{}{}

\section{Carry-Forward Gap Fill}

\boxnote{SoftRed}{Why this addendum is here.}{The May 1 plan itself asks for
\term{Urban environments} and \term{Resources}. However, the April 30 workbook
also asks you to carry forward one \term{Food and health} case detail and one
\term{Climate} case detail. Use this short addendum to stop those two case-bank
slots staying empty.}

\subsection{Food And Health: Haiti Cholera}

\rowcolors{2}{TableStripe}{white}
\begin{tabularx}{\textwidth}{L{0.18\textwidth}Y}
\toprule
\textbf{Case link} & \textbf{Exam-ready content} \\
\midrule
Place & Haiti, especially the 2010 cholera outbreak after the January 2010
earthquake \\
Issue & Disease diffusion, health inequality, unsafe water, sanitation, and
post-disaster vulnerability \\
Statistics & CDC: the 2010 outbreak was the worst in recent history, with
\term{over 820,000 cases} and \term{nearly 10,000 deaths}. CDC also notes the
January 2010 earthquake killed \term{over 200,000} people and displaced
\term{over 1 million}. \\
Causes & Earthquake damage, displacement, weak water and sanitation
infrastructure, poverty, crowded living conditions, and limited access to rapid
treatment \\
Impacts & Severe dehydration and death risk, pressure on health services,
uneven impacts on poorer and displaced communities, and national public-health
emergency \\
Response & Surveillance, laboratory capacity, cholera treatment centres, oral
rehydration, WASH measures, health-worker training, and oral cholera vaccine
campaigns \\
Evaluation & Emergency treatment and WASH responses reduced deaths, but the
case shows that disease control is limited if safe water, sanitation, housing,
and health-system capacity remain weak \\
\bottomrule
\end{tabularx}
\rowcolors{2}{}{}

\subsection{Climate: Bangladesh Cyclone Preparedness}

\rowcolors{2}{TableStripe}{white}
\begin{tabularx}{\textwidth}{L{0.18\textwidth}Y}
\toprule
\textbf{Case link} & \textbf{Exam-ready content} \\
\midrule
Place & Coastal Bangladesh, Bay of Bengal \\
Issue & Climate vulnerability and resilience to tropical cyclones, storm surge,
flooding, and coastal exposure \\
Statistics & World Bank: Cyclone Bhola in 1970 killed about \term{300,000}
people, while Cyclone Sidr in 2007 killed about \term{3,500}; Cyclone Fani in
2018 killed fewer than \term{10}, and Cyclone Amphan in 2020 killed around
\term{20}. Bangladesh Red Crescent: the Cyclone Preparedness Programme has
about \term{76,020 volunteers}. \\
Causes & Low-lying deltaic coast, high population density, storm surge exposure,
poverty, housing vulnerability, and climate-related hazard risk \\
Impacts & Deaths, displacement, livelihood damage, infrastructure loss, crop and
salinity impacts, and pressure on coastal communities \\
Response & Cyclone shelters, early warning, evacuation planning, volunteers,
coastal infrastructure, communication networks, and community preparedness \\
Evaluation & Bangladesh shows resilience can greatly reduce deaths, but risk
remains because exposure, poverty, housing quality, salinity, and future
climate change continue to threaten lives and livelihoods \\
\bottomrule
\end{tabularx}
\rowcolors{2}{}{}

\boxnote{SoftTint}{Use timing.}{Use Haiti before the May 2 Paper 1 mock if the
Food and health prompt involves disease, health inequality, water, sanitation,
or stakeholders. Use Bangladesh later for Paper 2 climate vulnerability and
resilience.}

\section{25-Minute Practice Block}

Use the workbook file
\term{2026-05-01-case-study-rapid-recall-workbook.pdf}. The aim is one rapid
recall sheet, not a new set of long notes.

\subsection{Minute-By-Minute Plan}

\rowcolors{2}{TableStripe}{white}
\begin{tabularx}{\textwidth}{L{0.15\textwidth}Y Y}
\toprule
\textbf{Time} & \textbf{Task} & \textbf{Output} \\
\midrule
0--4 min & Closed-book dump of all remembered case names and statistics & Raw
case list \\
4--8 min & Choose the best two Urban and best two Resources examples & Four
case rows \\
8--17 min & Fill the six-link chain for each case & Place, issue, causes,
impacts, response, evaluation \\
17--21 min & Add one statistic and one limitation to each case & Exam-ready
evidence \\
21--25 min & Test recall without looking, then mark uncertain facts & Error
log and check list \\
\bottomrule
\end{tabularx}
\rowcolors{2}{}{}

\subsection{One-Sentence Frames}

\boxnote{SoftGreen}{Urban frame.}{In \term{place}, \term{issue} developed
because \term{causes}, creating \term{impacts}; \term{response} tried to
manage it, but success was uneven because \term{evaluation}.}

\vspace{0.45em}

\boxnote{SoftBlue}{Resource frame.}{In \term{place}, \term{resource pressure}
was caused by \term{drivers}, affecting \term{people / environment};
\term{stewardship response} improved security, but it was limited by
\term{cost / equity / side effect / time scale}.}

\section{Source Notes}

Use these sources only to protect statistics. In the exam, write the case
chain, not a bibliography.

\begin{itemize}
  \item U.S. Census Bureau, 1950 census historical facts and Detroit
  QuickFacts: \href{https://www.census.gov/about/history/historical-censuses-and-surveys/decade-facts.1950.html}{1950 facts}
  and \href{https://www.census.gov/quickfacts/fact/table/detroitcitymichigan/POP645223}{Detroit QuickFacts}
  \item New Development Bank:
  \href{https://www.ndb.int/project/brazil-curitibas-bus-rapid-transit-rideability-improvement-project/}{Curitiba BRT Rideability Improvement Project}
  \item State Council of the People's Republic of China, 2018 solid waste
  import restrictions:
  \href{https://english.www.gov.cn/policies/policy_watch/2018/02/05/content_281476036805432.htm}{Tighter rules wielded on solid waste imports}
  \item ITU:
  \href{https://www.itu.int/en/ITU-D/Environment/Pages/Publications/The-Global-E-waste-Monitor-2024.aspx}{Global E-waste Monitor 2024}
  \item Singapore Ministry of Sustainability and the Environment, Water:
  \href{https://www.mse.gov.sg/policies/water/}{Water policy page}
  \item CDC:
  \href{https://www.cdc.gov/orr/responses/haiti-cholera-outbreak.html}{Haiti Cholera Outbreak}
  and
  \href{https://archive.cdc.gov/www_cdc_gov/cholera/haiti/2010-outbreak-response.html}{2010 Haiti Cholera Outbreak and CDC Response}
  \item World Bank:
  \href{https://www.worldbank.org/en/results/2022/08/29/shelter-from-the-storm-protecting-bangladesh-s-coastal-communities-from-natural-disasters}{Shelter from the Storm}
  and Bangladesh Red Crescent Society:
  \href{https://bdrcs.org/cyclone-preparedness-programm-cpp/}{Cyclone Preparedness Programme}
\end{itemize}

\section{End-Of-Day Output}

You are finished with May 1 if you can write these four chains from memory:

\begin{enumerate}
  \item Detroit: population loss $\rightarrow$ vacancy / fiscal stress
  $\rightarrow$ selective reinvestment $\rightarrow$ uneven regeneration
  $\rightarrow$ equity evaluation
  \item Curitiba: urban growth / congestion $\rightarrow$ BRT and corridor
  upgrades $\rightarrow$ mobility benefits $\rightarrow$ maintenance and
  ridership limits
  \item China e-waste: high consumption $\rightarrow$ waste trade
  $\rightarrow$ pollution and lost materials $\rightarrow$ import restrictions
  / circular economy $\rightarrow$ shifted responsibility
  \item Singapore water: physical scarcity and rising demand $\rightarrow$
  Four National Taps $\rightarrow$ more resilient supply $\rightarrow$ energy,
  cost, dependency, and demand-management limits
\end{enumerate}

\boxnote{SoftTint}{Carry forward.}{Saturday 2026-05-02 is the full Paper 1
simulation. Bring this case-study bank to the mock and use the examples only
where they answer the exact prompt.}

\end{document}
